\section{金的生产、法律与北方宗教社会}

\subsection{燕京接收的市场}

金取燕京并与宋交接。把它放在，可以看见商业、法律和宗教并非彼此分离的专题：货物与人员移动需要道路、渡口、票据、度量和裁判，寺院、学校或其他团体又可能保存捐施、契约、技术与身份的局部痕迹。名义移交不等于地方实控 这种联系不应被写成制度自然扩散，而应理解为不同机构在具体地点反复协调的结果。

这一判断的年序主要由。编年和本纪能够标出诏令、战争、港口或政权转换的时间，制度汇编能够说明某项规范的名称，地方文书、碑刻、船货、窑址和建筑遗存则保存更细小却更受限的证据。它们不能彼此替代：法律条文首先说明应当如何办理，诉讼或契约才可能让人看见一次具体实践；一处沉船说明一批货物的航程，不等于一种国家贸易政策。

因此，讨论市场不能只用税额或物价，讨论宗教也不能只用朝廷的禁令或褒奖。商人、船户、工匠、胥吏、僧道、学生、军人和农户接触制度的方式不同，且常在同一城市或乡村并存。现存材料特别偏向有书写能力、可立碑或被官府登记的人群；没有留下契约和墓志的人并非没有参与交换、劳役或信仰，只是较难被直接听见。

财政与司法的边界也并不清楚。征税、借贷、盐茶、关市、土地、寺产和工程往往都需要分类、登记、催纳和裁决。记录中的数额必须保留币种、实物单位、机关和地域，不能换算成一个跨政权的繁荣或压迫指数。反过来，一项地方个案也不足以否定制度文本。较可靠的写法，是让国家规范、地方执行和物质遗存相互限制，并明确何处仍只有单方叙述。

宗教与知识传播提供了另一条证据线。经卷、印本、寺院题记、书院文献、技术书和教育条目能够显示文本如何被抄写、雕印、施入或教学；它们不能直接证明所有居民读过同一文本，也不能将捐施者的愿望等同于公共信仰。语言、文字、官号和职业同样应分别对待，不可由其中之一推出稳定的族群或身份。

所以，金取燕京并与宋交接 的历史价值在于让：谁能够进入文书，谁承担运输或劳役，谁在市场、寺院、学校和衙门之间转换位置。其余尚不可见的部分，本书以来源限制标出，避免用连贯叙事填补材料的沉默。

\subsection{中原直辖与文书}

刘豫政权废除后直接经营。把它放在，可以看见商业、法律和宗教并非彼此分离的专题：货物与人员移动需要道路、渡口、票据、度量和裁判，寺院、学校或其他团体又可能保存捐施、契约、技术与身份的局部痕迹。政权转换有地区时间差 这种联系不应被写成制度自然扩散，而应理解为不同机构在具体地点反复协调的结果。

这一判断的年序主要由。编年和本纪能够标出诏令、战争、港口或政权转换的时间，制度汇编能够说明某项规范的名称，地方文书、碑刻、船货、窑址和建筑遗存则保存更细小却更受限的证据。它们不能彼此替代：法律条文首先说明应当如何办理，诉讼或契约才可能让人看见一次具体实践；一处沉船说明一批货物的航程，不等于一种国家贸易政策。

因此，讨论市场不能只用税额或物价，讨论宗教也不能只用朝廷的禁令或褒奖。商人、船户、工匠、胥吏、僧道、学生、军人和农户接触制度的方式不同，且常在同一城市或乡村并存。现存材料特别偏向有书写能力、可立碑或被官府登记的人群；没有留下契约和墓志的人并非没有参与交换、劳役或信仰，只是较难被直接听见。

财政与司法的边界也并不清楚。征税、借贷、盐茶、关市、土地、寺产和工程往往都需要分类、登记、催纳和裁决。记录中的数额必须保留币种、实物单位、机关和地域，不能换算成一个跨政权的繁荣或压迫指数。反过来，一项地方个案也不足以否定制度文本。较可靠的写法，是让国家规范、地方执行和物质遗存相互限制，并明确何处仍只有单方叙述。

宗教与知识传播提供了另一条证据线。经卷、印本、寺院题记、书院文献、技术书和教育条目能够显示文本如何被抄写、雕印、施入或教学；它们不能直接证明所有居民读过同一文本，也不能将捐施者的愿望等同于公共信仰。语言、文字、官号和职业同样应分别对待，不可由其中之一推出稳定的族群或身份。

所以，刘豫政权废除后直接经营 的历史价值在于让：谁能够进入文书，谁承担运输或劳役，谁在市场、寺院、学校和衙门之间转换位置。其余尚不可见的部分，本书以来源限制标出，避免用连贯叙事填补材料的沉默。

\subsection{中都营建与手工业}

金迁都中都。把它放在，可以看见商业、法律和宗教并非彼此分离的专题：货物与人员移动需要道路、渡口、票据、度量和裁判，寺院、学校或其他团体又可能保存捐施、契约、技术与身份的局部痕迹。遗址范围不等于全城生产 这种联系不应被写成制度自然扩散，而应理解为不同机构在具体地点反复协调的结果。

这一判断的年序主要由。编年和本纪能够标出诏令、战争、港口或政权转换的时间，制度汇编能够说明某项规范的名称，地方文书、碑刻、船货、窑址和建筑遗存则保存更细小却更受限的证据。它们不能彼此替代：法律条文首先说明应当如何办理，诉讼或契约才可能让人看见一次具体实践；一处沉船说明一批货物的航程，不等于一种国家贸易政策。

因此，讨论市场不能只用税额或物价，讨论宗教也不能只用朝廷的禁令或褒奖。商人、船户、工匠、胥吏、僧道、学生、军人和农户接触制度的方式不同，且常在同一城市或乡村并存。现存材料特别偏向有书写能力、可立碑或被官府登记的人群；没有留下契约和墓志的人并非没有参与交换、劳役或信仰，只是较难被直接听见。

财政与司法的边界也并不清楚。征税、借贷、盐茶、关市、土地、寺产和工程往往都需要分类、登记、催纳和裁决。记录中的数额必须保留币种、实物单位、机关和地域，不能换算成一个跨政权的繁荣或压迫指数。反过来，一项地方个案也不足以否定制度文本。较可靠的写法，是让国家规范、地方执行和物质遗存相互限制，并明确何处仍只有单方叙述。

宗教与知识传播提供了另一条证据线。经卷、印本、寺院题记、书院文献、技术书和教育条目能够显示文本如何被抄写、雕印、施入或教学；它们不能直接证明所有居民读过同一文本，也不能将捐施者的愿望等同于公共信仰。语言、文字、官号和职业同样应分别对待，不可由其中之一推出稳定的族群或身份。

所以，金迁都中都 的历史价值在于让：谁能够进入文书，谁承担运输或劳役，谁在市场、寺院、学校和衙门之间转换位置。其余尚不可见的部分，本书以来源限制标出，避免用连贯叙事填补材料的沉默。

\subsection{宋金关市与法令}

和议与再战交替。把它放在，可以看见商业、法律和宗教并非彼此分离的专题：货物与人员移动需要道路、渡口、票据、度量和裁判，寺院、学校或其他团体又可能保存捐施、契约、技术与身份的局部痕迹。条约不能自动说明地方交易 这种联系不应被写成制度自然扩散，而应理解为不同机构在具体地点反复协调的结果。

这一判断的年序主要由。编年和本纪能够标出诏令、战争、港口或政权转换的时间，制度汇编能够说明某项规范的名称，地方文书、碑刻、船货、窑址和建筑遗存则保存更细小却更受限的证据。它们不能彼此替代：法律条文首先说明应当如何办理，诉讼或契约才可能让人看见一次具体实践；一处沉船说明一批货物的航程，不等于一种国家贸易政策。

因此，讨论市场不能只用税额或物价，讨论宗教也不能只用朝廷的禁令或褒奖。商人、船户、工匠、胥吏、僧道、学生、军人和农户接触制度的方式不同，且常在同一城市或乡村并存。现存材料特别偏向有书写能力、可立碑或被官府登记的人群；没有留下契约和墓志的人并非没有参与交换、劳役或信仰，只是较难被直接听见。

财政与司法的边界也并不清楚。征税、借贷、盐茶、关市、土地、寺产和工程往往都需要分类、登记、催纳和裁决。记录中的数额必须保留币种、实物单位、机关和地域，不能换算成一个跨政权的繁荣或压迫指数。反过来，一项地方个案也不足以否定制度文本。较可靠的写法，是让国家规范、地方执行和物质遗存相互限制，并明确何处仍只有单方叙述。

宗教与知识传播提供了另一条证据线。经卷、印本、寺院题记、书院文献、技术书和教育条目能够显示文本如何被抄写、雕印、施入或教学；它们不能直接证明所有居民读过同一文本，也不能将捐施者的愿望等同于公共信仰。语言、文字、官号和职业同样应分别对待，不可由其中之一推出稳定的族群或身份。

所以，和议与再战交替 的历史价值在于让：谁能够进入文书，谁承担运输或劳役，谁在市场、寺院、学校和衙门之间转换位置。其余尚不可见的部分，本书以来源限制标出，避免用连贯叙事填补材料的沉默。

\subsection{寺院碑志与家族网络}

金代碑志寺院材料延续。把它放在，可以看见商业、法律和宗教并非彼此分离的专题：货物与人员移动需要道路、渡口、票据、度量和裁判，寺院、学校或其他团体又可能保存捐施、契约、技术与身份的局部痕迹。追赠官衔不可当作实际职务 这种联系不应被写成制度自然扩散，而应理解为不同机构在具体地点反复协调的结果。

这一判断的年序主要由。编年和本纪能够标出诏令、战争、港口或政权转换的时间，制度汇编能够说明某项规范的名称，地方文书、碑刻、船货、窑址和建筑遗存则保存更细小却更受限的证据。它们不能彼此替代：法律条文首先说明应当如何办理，诉讼或契约才可能让人看见一次具体实践；一处沉船说明一批货物的航程，不等于一种国家贸易政策。

因此，讨论市场不能只用税额或物价，讨论宗教也不能只用朝廷的禁令或褒奖。商人、船户、工匠、胥吏、僧道、学生、军人和农户接触制度的方式不同，且常在同一城市或乡村并存。现存材料特别偏向有书写能力、可立碑或被官府登记的人群；没有留下契约和墓志的人并非没有参与交换、劳役或信仰，只是较难被直接听见。

财政与司法的边界也并不清楚。征税、借贷、盐茶、关市、土地、寺产和工程往往都需要分类、登记、催纳和裁决。记录中的数额必须保留币种、实物单位、机关和地域，不能换算成一个跨政权的繁荣或压迫指数。反过来，一项地方个案也不足以否定制度文本。较可靠的写法，是让国家规范、地方执行和物质遗存相互限制，并明确何处仍只有单方叙述。

宗教与知识传播提供了另一条证据线。经卷、印本、寺院题记、书院文献、技术书和教育条目能够显示文本如何被抄写、雕印、施入或教学；它们不能直接证明所有居民读过同一文本，也不能将捐施者的愿望等同于公共信仰。语言、文字、官号和职业同样应分别对待，不可由其中之一推出稳定的族群或身份。

所以，金代碑志寺院材料延续 的历史价值在于让：谁能够进入文书，谁承担运输或劳役，谁在市场、寺院、学校和衙门之间转换位置。其余尚不可见的部分，本书以来源限制标出，避免用连贯叙事填补材料的沉默。

\subsection{证据登记与地方实践的距离}

商业、海域、法律与宗教制度能够相互照明，却不能彼此代证。市场税、盐茶、船货、寺产、田宅、学额或僧道名籍的数字，都产生于某个抄录者、某个机关和某次治理需要；同一数字离开币种、单位、时间和地域，便不再是可比较的事实。阅读它们时，应先问记录者为何要把某一项写成收入、欠额、功德、违禁或灾异，再问这一分类是否被地方实际接受。制度反复申明本身常常表明执行存在摩擦，而不是证明命令无所不至。

地方材料也有相反的诱惑。一处港口、一份契约、一通碑刻或一批窖藏，能够让某些人名、货物、路线和纠纷变得具体，却仍是偶然保存下来的证据。它不能自动说明邻近地区，更不能代表所有身份群体。将制度文本同文书、物质遗存和对方叙事放在同一问题中，并非要求它们给出完全相同的答案；恰恰是它们的不一致，使我们能够辨认国家分类、地方利益和宗教或市场网络各自运行的边界。

对多政权时代尤其如此。国号、族名、官号和宗教语言常在外交、司法和记功的场景中使用，必须保留其原有语境。以一种现代的、固定的身份替代当时复杂的地域、职业、书写与法律关系，会使商业流动和地方协商重新变得不可见。因而本书只在材料允许的范围内描述人群与机构；没有证据支持的动机、规模和普遍性，宁可留作待查问题。

